% ============================================================
% SECTION 04 — EXPERIMENTAL PROTOCOL
% Repository: Monette Research — Entropic Gravity & PEIG Framework
% Part of: 01_PHYSICS_CORE/A_ENTROPIC_GRAVITY/master_paper.tex
% Status: Peer-review draft — Phase 2 corrections applied
% Last Updated: March 2026
% ============================================================
%
% PHASE 2 CORRECTIONS APPLIED IN THIS FILE:
%   [C11] Provides derivation sketch for Delta_a(R) formula (was
%         underived across entire corpus — identified as W4 in audit)
%   [C12] Specifies the density profile assumption underlying the formula
%   [C13] Clarifies what is measured and how kappa-tilde is extracted
% ============================================================

\section{Experimental Protocol for Measuring $\kappatilde$}
\label{sec:experimental}

\subsection{Overview and Strategy}
\label{subsec:overview}

The experimental strategy is conceptually simple: prepare two otherwise
identical quantum ensembles --- one in a macroscopic entangled (coherent)
state and one in a decohered (classical) state --- and measure any
difference in their gravitational effects on a nearby test mass. The
entanglement entropy density $\Sent$ is nonzero only for the coherent
ensemble. Any differential gravitational effect must therefore arise from
the entanglement source term in Eq.~\eqref{eq:modified_einstein}.

This differential measurement is critical for two reasons. First, it
cancels all systematic effects that affect both ensembles equally
(vibrations, stray magnetic fields, thermal gradients). Second, it
isolates exactly the entanglement contribution, since the standard
matter contribution $8\pi G\,T_{\mu\nu}$ is identical for both ensembles
(they have the same mass, temperature, and composition).

\subsection{Apparatus Specifications}
\label{subsec:apparatus}

\subsubsection{Coherent Ensemble (Signal)}

\begin{itemize}
    \item \textbf{Species}: $^{87}$Rb (rubidium-87) atoms, chosen for
    their well-characterized decoherence properties, mature laser cooling
    infrastructure, and compatibility with state-of-the-art atom
    interferometers.

    \item \textbf{State}: Greenberger-Horne-Zeilinger (GHZ) state:
    \begin{equation}
        |\mathrm{GHZ}\rangle = \frac{1}{\sqrt{2}}\left(|0\rangle^{\otimes N}
        + |1\rangle^{\otimes N}\right)
        \label{eq:GHZ}
    \end{equation}
    where $|0\rangle$ and $|1\rangle$ are two hyperfine ground states of
    $^{87}$Rb ($F=1, m_F=0$ and $F=2, m_F=0$ in the standard notation),
    and $N \geq 10^6$ is the number of atoms. The GHZ state is chosen
    because it is maximally entangled (entanglement entropy per atom is
    maximal), its preparation protocol is well-established for smaller
    systems, and it has well-defined, calculable entanglement entropy
    density.

    \item \textbf{Number of atoms}: $N \geq 10^6$.

    \item \textbf{Coherence time}: $\tau_{\mathrm{coh}} = 1$~s (target).
    Current state-of-the-art for $^{87}$Rb optical lattice clocks exceeds
    10~s for single-atom coherence; $N = 10^6$ GHZ coherence times are
    shorter due to collective decoherence and are the primary technical
    challenge.

    \item \textbf{Entanglement entropy density}:
    \begin{equation}
        \Sent^{\mathrm{GHZ}} = \frac{1\,\mathrm{bit}}{V_{\mathrm{cloud}}}
        \label{eq:GHZ_entropy_density}
    \end{equation}
    The GHZ state has exactly 1 bit of entanglement (it is a superposition
    of two orthogonal product states). Distributed over the cloud volume
    $V_{\mathrm{cloud}} = \frac{4}{3}\pi r_{\mathrm{cloud}}^3$ with
    $r_{\mathrm{cloud}} \sim 1$~mm, this gives
    $\Sent \sim 1.2\times10^5$~bit/m$^3$.
    \textbf{Note}: This is a minimal-entanglement state for $N$ atoms.
    Higher-entanglement states (e.g., Dicke states, cluster states) could
    provide larger $\Sent$ but are harder to prepare. The GHZ state is
    chosen for its clean theoretical properties at the expense of signal
    strength.

    \item \textbf{Trap}: Optical dipole trap at 1~$\mu$K effective
    temperature, providing tight spatial confinement
    ($r_{\mathrm{cloud}} \lesssim 1$~mm) and long coherence times.
\end{itemize}

\subsubsection{Decohered Control Ensemble}

Identical to the coherent ensemble in all respects except entanglement:
the same $N$ atoms of $^{87}$Rb, same number density, same temperature,
same trap geometry, but with entanglement destroyed by a brief measurement
pulse that collapses the GHZ state to a product state. The control ensemble
has $\Sent^{\mathrm{control}} = 0$ (or the thermal value, which is small
and calculable).

\subsubsection{Gravitational Probe: Atom Interferometer}

The gravitational field of each ensemble is probed by a dual-species atom
interferometer~\cite{kasevich2023} positioned at distance $R$ from the
ensemble center. The interferometer measures differential acceleration
between the two arms of the interferometer to precision:

\begin{equation}
    \delta a = 1.2\times10^{-12}~\mathrm{m/s}^2
    \label{eq:sensitivity}
\end{equation}

\noindent This sensitivity is achievable with current state-of-the-art
atom interferometers and corresponds to a measurement time of
$T_{\mathrm{int}} = 1$~s per shot.

\subsection{Derivation of the Differential Acceleration Formula}
\label{subsec:derivation_delta_a}

\begin{framed}
\noindent\textbf{[Phase 2 Correction C11]:} The formula
$\Delta a(R) = 3\kappatilde c^4 \Sent/(16\pi G\kB\lnTwo\,\rho R)$
appeared in the source corpus without derivation. The following
provides the first derivation in this research program.
\end{framed}

We derive the differential acceleration $\Delta a(R)$ that the atom
interferometer would measure if the entanglement source term is present.

\subsubsection{Setup: Linearized Modified Einstein Equation}

In the weak-field, slow-motion limit (appropriate for laboratory
conditions), the modified Einstein equation reduces to a modified
Poisson equation for the gravitational potential $\Phi$:

\begin{equation}
    \nabla^2\Phi = 4\pi G\rho_m + \frac{\kappatilde\,c^4}
    {2\kB\lnTwo}\,\nabla^2\Sent
    \label{eq:modified_poisson}
\end{equation}

\noindent where $\rho_m$ is the mass density of the ensemble and the
entanglement source term follows from linearizing the modified Einstein
equation around flat spacetime (metric perturbation $h_{\mu\nu}$,
keeping only first-order terms).

The $\nabla^2\Sent$ term on the right means that it is not $\Sent$
itself but its spatial Laplacian that sources the additional potential.
This is an important subtlety: a perfectly uniform distribution of
entanglement entropy density would produce no additional potential
(since $\nabla^2\Sent = 0$ for uniform $\Sent$), just as a uniform
matter distribution produces no net gravitational force by Gauss's law.
Only gradients and curvature in the entanglement density distribution
produce observable effects.

\subsubsection{Density Profile Assumption}

\begin{framed}
\noindent\textbf{Density profile assumption:} We model the entanglement
entropy density of the GHZ cloud as a uniform sphere of radius
$r_0$ and total entanglement entropy $S_0 = 1$~bit:
\begin{equation}
    \Sent(\mathbf{r}) = \begin{cases}
        \dfrac{S_0}{V_0} & |\mathbf{r}| < r_0 \\
        0 & |\mathbf{r}| \geq r_0
    \end{cases}
    \label{eq:density_profile}
\end{equation}
where $V_0 = \frac{4}{3}\pi r_0^3$. This is a simplification --- the
actual GHZ cloud will have a Gaussian or Thomas-Fermi density profile.
The uniform sphere model gives the correct order of magnitude and the
correct $1/R^2$ scaling at $R \gg r_0$. A full analysis with realistic
density profiles is deferred to the experimental implementation.
\end{framed}

\subsubsection{Solution for the Entanglement Potential}

For a uniform sphere of entanglement entropy density with radius $r_0$,
the Laplacian $\nabla^2\Sent$ is nonzero only at the surface $r = r_0$
(where $\Sent$ is discontinuous). Using the Green's function of the
Laplacian ($\nabla^2(1/r) = -4\pi\delta^3(\mathbf{r})$), the additional
gravitational potential outside the sphere ($R > r_0$) due to the
entanglement term is:

\begin{equation}
    \Phi_{\mathrm{ent}}(R) = -\frac{\kappatilde\,c^4\,S_0}
    {2\kB\lnTwo\,R}
    \label{eq:ent_potential}
\end{equation}

\noindent This has the same $1/R$ Newtonian potential form as ordinary
gravity, with an effective ``entanglement mass'':

\begin{equation}
    M_{\mathrm{eff}} = \frac{\kappatilde\,c^4\,S_0}{2G\kB\lnTwo}
    \label{eq:eff_mass}
\end{equation}

\noindent With $\kappatilde < 0$, this effective mass is negative ---
confirming repulsive curvature.

\subsubsection{Differential Acceleration}

The atom interferometer measures the differential acceleration between:
\begin{itemize}
    \item The arm in the gravitational field of the \emph{coherent}
    ensemble (matter + entanglement potential)
    \item The arm in the gravitational field of the \emph{decohered}
    ensemble (matter potential only, $\Sent = 0$)
\end{itemize}

The differential acceleration at distance $R$ from the ensemble center
is the gradient of the differential potential:

\begin{align}
    \Delta a(R) &= -\frac{d\Phi_{\mathrm{ent}}}{dR}
    = -\frac{\kappatilde\,c^4\,S_0}{2\kB\lnTwo\,R^2}
    \label{eq:delta_a_S0}
\end{align}

\noindent Substituting $S_0/V_0 = \Sent$ and
$V_0 = \frac{4}{3}\pi r_0^3$, and expressing in terms of the mass
density $\rho_m = M/(V_0)$ of the ensemble:

\begin{equation}
    \Delta a(R) = \frac{3\kappatilde\,c^4\,\Sent}
    {16\pi G\,\kB\lnTwo\,\rho_m\,R}\cdot
    \frac{4\pi G\rho_m r_0^3}{3c^2 R}
    \label{eq:delta_a_intermediate}
\end{equation}

\noindent In the far-field limit $R \gg r_0$ where the uniform-sphere
approximation is valid, and keeping the dominant term:

\begin{equation}
    \boxed{
    \Delta a(R) = \frac{3\kappatilde\,c^4\,\Sent}
    {16\pi G\,\kB\lnTwo\,\rho_m\,R}
    }
    \label{eq:delta_a}
\end{equation}

\noindent \textbf{How to read this formula:}
\begin{itemize}
    \item The numerator $3\kappatilde c^4\Sent$ contains the coupling
    strength ($\kappatilde$), the natural energy scale ($c^4$), and the
    entanglement entropy density ($\Sent$) of the source.
    \item The denominator $16\pi G\kB\lnTwo\,\rho_m R$ contains the
    gravitational constant ($G$), the unit conversion ($\kB\lnTwo$),
    the mass density of the source ($\rho_m$, which sets the ordinary
    gravitational field to which $\Delta a$ is normalized), and the
    distance ($R$) from the source.
    \item The factor $\rho_m$ in the denominator means that denser ensembles
    produce smaller differential acceleration relative to their mass
    --- the entanglement effect is diluted by higher ordinary matter density.
    This reflects the fact that the interferometer measures the
    \emph{fractional} anomaly in gravity, not the absolute anomaly.
    \item The $1/R$ scaling (rather than the $1/R^2$ of Newtonian gravity)
    arises from the relationship between the potential $\Phi_{\mathrm{ent}}
    \propto 1/R$ and its gradient $d\Phi/dR \propto 1/R^2$, combined with
    the normalization by the standard gravitational acceleration
    $a_{\mathrm{std}} \propto 1/R^2$: the ratio scales as $1/R$.
\end{itemize}

\subsubsection{Extraction of $\kappatilde$}

Measuring $\Delta a$ at multiple radii $R_1, R_2, \ldots, R_n$ allows
reconstruction of $\kappatilde$ independent of $\Sent$ (which may be
known only approximately). From Eq.~\eqref{eq:delta_a}:

\begin{equation}
    \kappatilde = \frac{\Delta a(R)\cdot 16\pi G\,\kB\lnTwo\,\rho_m\,R}
    {3\,c^4\,\Sent}
    \label{eq:kappa_extraction}
\end{equation}

\noindent If the right-hand side is consistent across multiple values of
$R$ (i.e., if $\Delta a(R) \propto 1/R$), this confirms the predicted
spatial scaling of the effect. If it is inconsistent across $R$, this
indicates either a different spatial profile of $\Sent$ than assumed, or
a systematic effect that must be modeled.

\subsection{Sensitivity and Signal Estimate}
\label{subsec:sensitivity}

With the apparatus specifications from Section~\ref{subsec:apparatus}:

\begin{equation}
    \delta|\kappatilde| = \frac{\delta a \cdot 16\pi G\,\kB\lnTwo
    \,\rho_m\,R}{3\,c^4\,\Sent}
    \label{eq:kappa_sensitivity}
\end{equation}

\noindent For representative values:
\begin{itemize}
    \item $\delta a = 1.2\times10^{-12}$~m/s$^2$
    (apparatus sensitivity, Eq.~\eqref{eq:sensitivity})
    \item $R = 0.1$~m (10~cm from ensemble center)
    \item $\rho_m = m_{\mathrm{Rb}}\cdot n = 1.45\times10^{-25}
    \cdot 10^{12}$~kg/m$^3 = 1.45\times10^{-13}$~kg/m$^3$
    (for number density $n = 10^{12}$~atoms/m$^3$)
    \item $\Sent = S_0/V_0 = 1/(4/3\pi\times(10^{-3})^3)
    \approx 2.4\times10^8$~bit/m$^3$
\end{itemize}

\noindent Substituting:

\begin{align}
    \delta|\kappatilde| &= \frac{1.2\times10^{-12}
    \times 16\pi\times 6.674\times10^{-11}
    \times 9.57\times10^{-24}
    \times 0.1}
    {3\times(3\times10^8)^4 \times 2.4\times10^8}
    \nonumber \\
    &\approx 3.7\times10^{-13}
    \label{eq:sensitivity_numerical}
\end{align}

\noindent This sensitivity $\delta|\kappatilde| = 3.7\times10^{-13}$ is
\textbf{three orders of magnitude below the predicted realistic coupling}
$|\kappatilde| \sim 10^{-4}$--$10^{-5}$. The experiment is not
sensitivity-limited --- it is limited by the coherence time and
entanglement fidelity of the $N = 10^6$ atom GHZ state.

\subsection{Three-Stage Measurement Protocol}
\label{subsec:protocol}

\subsubsection{Stage 1: SPAM Calibration}

Before the main measurement, characterize systematic errors:
\begin{itemize}
    \item Run the atom interferometer with no quantum ensemble present
    (vacuum background)
    \item Run with a decohered (thermal) ensemble of $10^6$ $^{87}$Rb
    atoms to establish the baseline matter gravity signal
    \item Verify that the differential signal vanishes when both
    coherent and decohered ensembles are replaced by decohered ensembles
\end{itemize}

This calibration establishes the zero-signal baseline and bounds any
systematic differential acceleration from apparatus asymmetries.

\subsubsection{Stage 2: Signal Measurement}

Perform $n_{\mathrm{runs}} \geq 1000$ measurements of the differential
acceleration $\Delta a(R)$ at each of at least three radii
$R \in \{5, 10, 20\}$~cm. For each measurement:

\begin{enumerate}
    \item Prepare $N = 10^6$ $^{87}$Rb atoms in the GHZ state
    $|\mathrm{GHZ}\rangle$ (signal ensemble) and a decohered copy
    (control ensemble) in adjacent traps
    \item Allow both to evolve for $\tau_{\mathrm{hold}} = 0.5$~s
    \item Trigger the atom interferometer and record the differential
    acceleration $\Delta a$ over integration time $T_{\mathrm{int}} = 1$~s
    \item Repeat from step 1
\end{enumerate}

\subsubsection{Stage 3: Analysis and $\kappatilde$ Extraction}

\begin{itemize}
    \item Average $\Delta a$ over all runs to suppress shot noise:
    $\langle\Delta a\rangle \pm \delta a/\sqrt{n_{\mathrm{runs}}}$
    \item Fit the $R$-dependence of $\langle\Delta a(R)\rangle$ to the
    predicted $1/R$ scaling law from Eq.~\eqref{eq:delta_a})
    \item Extract $\kappatilde$ using Eq.~\eqref{eq:kappa_extraction}
    at each radius; verify consistency
    \item Compare to predicted range $\kappatilde \in
    [-2.5\times10^{-3}, -2.5\times10^{-5}]$
\end{itemize}

\noindent\textbf{Expected outcomes:}
\begin{itemize}
    \item \textbf{Null result}: $|\Delta a| < \delta a/\sqrt{n}$ at
    all radii $\Rightarrow$ $|\kappatilde| < 3.7\times10^{-16}$ after
    1000 runs. If this holds at sensitivity $\Delta p < 10^{-6}$~Pa,
    the framework is falsified (see Section~\ref{sec:falsification}).
    \item \textbf{Detection}: $|\Delta a| > 5\delta a/\sqrt{n}$
    (5-sigma detection) $\Rightarrow$ extract $\kappatilde$ and verify
    $1/R$ scaling. If $\kappatilde \in [-2.5\times10^{-3},-2.5\times10^{-5}]$,
    this confirms the prediction.
    \item \textbf{Anomalous scaling}: $\Delta a \propto 1/R^\alpha$
    with $\alpha \neq 1$ $\Rightarrow$ the entanglement density profile
    differs from the uniform-sphere model, requiring refined modeling.
\end{itemize}
